% GENERATED FILE. Edit canonical lessons, then run npm run generate:books.
% canonical-source-hash: fnv1a64:ca1a7763e4f1bb32
% canonical-lessons: BN-C211-chechano, BN-C211-rag-kora, BN-C211-lojja-paoya, BN-C211-jhogra-kora, BN-C211-joriye-dhora

\chapter{To shout, To get angry, To feel shy, To quarrel, To hug}
\label{ch:bn-to-shout-to-get-angry-to-feel-shy-to-quarrel-to-hug}
\hlchaptermodality{211}

\begin{chapteropening}
\textbf{By the end of this chapter:} \emph{I can say to shout, to get angry, to feel shy, to quarrel and to hug.}

Complete the last lesson of chapter 211: I can say to shout, to get angry, to feel shy, to quarrel and to hug.
\end{chapteropening}

\section[\texorpdfstring{\bn{চেঁচানো} (chẽchāno)}{চেঁচানো (chẽchāno)}]{\bn{চেঁচানো} (chẽchāno) — to shout}

\label{lesson:BN-C211-chechano}



\begin{quote}
Before the new one: say the Bengali for to marry, then the Bengali for to survive.
\end{quote}

\subsection*{You'll want to know: \bn{চেঁচানো}}
\textbf{\bn{চেঁচানো}} — \emph{chẽchāno} — \textquotedblleft{}to shout\textquotedblright{}.

\begin{grammarlens}[title={What you've built}]
One more word for this chapter.
\end{grammarlens}

\subsection*{Your turn}
\begin{itemize}
\raggedright
  \item \emph{Say it:} \emph{chẽchāno}
  \item \emph{Say it:} \emph{chẽchāno}, once more
  \item \emph{Say it:} say \emph{chẽchāno}
  \item \emph{Recall:} say the Bengali for to marry, then the Bengali for to survive, then say \emph{chẽchāno} again
\end{itemize}

\begin{tcolorbox}[breakable,colback=teal!4,colframe=teal!35!black,title={Before you move on}]
What does \bn{চেঁচানো} mean? (\textquotedblleft{}To shout\textquotedblright{}.) Say it once more.
\end{tcolorbox}
\section[\texorpdfstring{\bn{রাগ} \bn{করা} (rāg kôrā)}{রাগ করা (rāg kôrā)}]{\bn{রাগ} \bn{করা} (rāg kôrā) — to get angry}

\label{lesson:BN-C211-rag-kora}



\begin{quote}
Before the new one: say the Bengali for to survive, then the Bengali for to shout.
\end{quote}

\subsection*{You'll want to know: \bn{রাগ} \bn{করা}}
\textbf{\bn{রাগ} \bn{করা}} — \emph{rāg kôrā} — \textquotedblleft{}to get angry\textquotedblright{}.

\begin{grammarlens}[title={What you've built}]
One more word for this chapter.
\end{grammarlens}

\subsection*{Your turn}
\begin{itemize}
\raggedright
  \item \emph{Say it:} \emph{rāg kôrā}
  \item \emph{Say it:} \emph{rāg kôrā}, once more
  \item \emph{Say it:} say \emph{rāg kôrā}
  \item \emph{Recall:} say the Bengali for to survive, then the Bengali for to shout, then say \emph{rāg kôrā} again
\end{itemize}

\begin{tcolorbox}[breakable,colback=teal!4,colframe=teal!35!black,title={Before you move on}]
What does \bn{রাগ} \bn{করা} mean? (\textquotedblleft{}To get angry\textquotedblright{}.) Say it once more.
\end{tcolorbox}
\section[\texorpdfstring{\bn{লজ্জা} \bn{পাওয়া} (lôjjā pāoyā)}{লজ্জা পাওয়া (lôjjā pāoyā)}]{\bn{লজ্জা} \bn{পাওয়া} (lôjjā pāoyā) — to feel shy}

\label{lesson:BN-C211-lojja-paoya}



\begin{quote}
Before the new one: say the Bengali for to shout, then the Bengali for to get angry.
\end{quote}

\subsection*{You'll want to know: \bn{লজ্জা} \bn{পাওয়া}}
\textbf{\bn{লজ্জা} \bn{পাওয়া}} — \emph{lôjjā pāoyā} — \textquotedblleft{}to feel shy\textquotedblright{}.

\begin{grammarlens}[title={What you've built}]
One more word for this chapter.
\end{grammarlens}

\subsection*{Your turn}
\begin{itemize}
\raggedright
  \item \emph{Say it:} \emph{lôjjā pāoyā}
  \item \emph{Say it:} \emph{lôjjā pāoyā}, once more
  \item \emph{Say it:} say \emph{lôjjā pāoyā}
  \item \emph{Recall:} say the Bengali for to shout, then the Bengali for to get angry, then say \emph{lôjjā pāoyā} again
\end{itemize}

\begin{tcolorbox}[breakable,colback=teal!4,colframe=teal!35!black,title={Before you move on}]
What does \bn{লজ্জা} \bn{পাওয়া} mean? (\textquotedblleft{}To feel shy\textquotedblright{}.) Say it once more.
\end{tcolorbox}
\section[\texorpdfstring{\bn{ঝগড়া} \bn{করা} (jhôgṛā kôrā)}{ঝগড়া করা (jhôgṛā kôrā)}]{\bn{ঝগড়া} \bn{করা} (jhôgṛā kôrā) — to quarrel}

\label{lesson:BN-C211-jhogra-kora}



\begin{quote}
Before the new one: say the Bengali for to get angry, then the Bengali for to feel shy.
\end{quote}

\subsection*{You'll want to know: \bn{ঝগড়া} \bn{করা}}
\textbf{\bn{ঝগড়া} \bn{করা}} — \emph{jhôgṛā kôrā} — \textquotedblleft{}to quarrel\textquotedblright{}.

\begin{grammarlens}[title={What you've built}]
One more word for this chapter.
\end{grammarlens}

\subsection*{Your turn}
\begin{itemize}
\raggedright
  \item \emph{Say it:} \emph{jhôgṛā kôrā}
  \item \emph{Say it:} \emph{jhôgṛā kôrā}, once more
  \item \emph{Say it:} say \emph{jhôgṛā kôrā}
  \item \emph{Recall:} say the Bengali for to get angry, then the Bengali for to feel shy, then say \emph{jhôgṛā kôrā} again
\end{itemize}

\begin{tcolorbox}[breakable,colback=teal!4,colframe=teal!35!black,title={Before you move on}]
What does \bn{ঝগড়া} \bn{করা} mean? (\textquotedblleft{}To quarrel\textquotedblright{}.) Say it once more.
\end{tcolorbox}
\section[\texorpdfstring{\bn{জড়িয়ে} \bn{ধরা} (jôṛiye dhôrā)}{জড়িয়ে ধরা (jôṛiye dhôrā)}]{\bn{জড়িয়ে} \bn{ধরা} (jôṛiye dhôrā) — to hug}

\label{lesson:BN-C211-joriye-dhora}



\begin{quote}
Before the new one: say the Bengali for to feel shy, then the Bengali for to quarrel.
\end{quote}

\subsection*{You'll want to know: \bn{জড়িয়ে} \bn{ধরা}}
\textbf{\bn{জড়িয়ে} \bn{ধরা}} — \emph{jôṛiye dhôrā} — \textquotedblleft{}to hug\textquotedblright{}.

\begin{grammarlens}[title={What you've built}]
One more word for this chapter.
\end{grammarlens}

\subsection*{Your turn}
\begin{itemize}
\raggedright
  \item \emph{Say it:} \emph{jôṛiye dhôrā}
  \item \emph{Say it:} \emph{jôṛiye dhôrā}, once more
  \item \emph{Say it:} say \emph{jôṛiye dhôrā}
  \item \emph{Recall:} say the Bengali for to feel shy, then the Bengali for to quarrel, then say \emph{jôṛiye dhôrā} again
\end{itemize}

\begin{tcolorbox}[breakable,colback=teal!4,colframe=teal!35!black,title={Before you move on}]
What does \bn{জড়িয়ে} \bn{ধরা} mean? (\textquotedblleft{}To hug\textquotedblright{}.) Say it once more.
\end{tcolorbox}
